
\documentclass{article}
\usepackage{colortbl}
\usepackage{makecell}
\usepackage{multirow}
\usepackage{supertabular}

\begin{document}

\newcounter{utterance}

\centering \large Interaction Transcript for game `wordle\_withcritic\_noexp', experiment `medium\_frequency\_words\_clue\_with\_critic', episode 5 with reflexion{-}Qwen3.5{-}4B.
\vspace{24pt}

{ \footnotesize  \setcounter{utterance}{1}
\setlength{\tabcolsep}{0pt}
\begin{supertabular}{c@{$\;$}|p{.15\linewidth}@{}p{.15\linewidth}p{.15\linewidth}p{.15\linewidth}p{.15\linewidth}p{.15\linewidth}}
    \# & $\;$A & \multicolumn{4}{c}{Game Master} & $\;\:$B\\
    \hline

    \theutterance \stepcounter{utterance}  
    & & \multicolumn{4}{p{0.6\linewidth}}{
        \cellcolor[rgb]{0.9,0.9,0.9}{
            \makecell[{{p{\linewidth}}}]{
                \texttt{\tiny{[P1$\langle$GM]}}
                \texttt{You are a language wizard who likes to guess words by using the given rules.} \\
\\ 
\texttt{Welcome to Wordle! You have six attempts to guess the target word, a valid English word of five lowercase letters (a{-}z). Please use the tag "guess:" to provide your guess.} \\
\\ 
\texttt{To help you make an informed guess, you will receive a clue for the word, such as} \\
\texttt{clue: "snowy white"} \\
\\ 
\texttt{Here is an example guess based on the clue: } \\
\texttt{guess: apple} \\
\\ 
\texttt{I will then indicate whether I agree or disagree with your guess, but agreeing with a guess does not confirm its correctness. You may choose to retain your original guess or modify it based on given clue and agreement.} \\
\\ 
\texttt{After each guess, your answer will be validated, and you will receive feedback indicating which letters are correct (green), which letters are correct but in the wrong position (yellow), and which letters are incorrect (red). This feedback can be useful in determining which letters to include or exclude in your next guess.} \\
\\ 
\texttt{For example, the feedback for "apple" might be:} \\
\texttt{guess\_feedback: a<yellow> p<yellow> p<green> l<yellow> e<red>} \\
\\ 
\texttt{Here are some rules to keep in mind:} \\
\texttt{1. Please do not generate any guess feedback or clue on your own} \\
\texttt{2. Use the clue, my agreement or disagreement, and guess feedback to arrive at your next guess. } \\
\\ 
\texttt{Do not generate any other text.} \\
\\ 
\texttt{Let's begin with your first guess.} \\
\\ 
\texttt{clue: tick off} \\
            }
        }
    }
    & \\ \\

    \theutterance \stepcounter{utterance}  
    & \multicolumn{4}{p{0.6\linewidth}}{
        \cellcolor[rgb]{0.9,0.9,0.9}{
            \makecell[{{p{\linewidth}}}]{
                \texttt{\tiny{[P1$\rangle$GM]}}
                \texttt{guess: tick} \\
            }
        }
    }
    & & \\ \\

    \theutterance \stepcounter{utterance}  
    & & & \multicolumn{2}{p{0.3\linewidth}}{
        \cellcolor[rgb]{0.9,0.9,0.9}{
            \makecell[{{p{\linewidth}}}]{
                \texttt{\tiny{[GM$|$GM]}}
                \texttt{Error: The length of the guessed word is not 5.} \\
            }
        }
    }
    & & \\ \\

    \theutterance \stepcounter{utterance}  
    & & & \multicolumn{2}{p{0.3\linewidth}}{
        \cellcolor[rgb]{0.9,0.9,0.9}{
            \makecell[{{p{\linewidth}}}]{
                \texttt{\tiny{[GM$|$GM]}}
                \texttt{game\_result = ABORT} \\
            }
        }
    }
    & & \\ \\

    \theutterance \stepcounter{utterance}  
    & \multicolumn{4}{p{0.6\linewidth}}{
        \cellcolor[rgb]{0.95,0.95,0.95}{
            \makecell[{{p{\linewidth}}}]{
                \texttt{\tiny{[P2$\rightarrow$P1]}}
                \texttt{Post{-}episode Reflexion (not sent to GM):} \\
\texttt{New plan: The last attempt was aborted because the response included the label "guess:" before the word, but the instructions explicitly stated to use the tag "guess:" to provide the guess, implying the format should be just the word after the tag or the tag was part of the instruction for the AI to output the guess itself. However, looking closely at the prompt: "Use the tag 'guess:' to provide your guess." and the example shows "guess: apple". The error in the previous attempt was likely that I output "guess: tick" but the system rejected it, possibly because the prompt said "Do not generate any other text" and I included the tag in a way that violated a hidden constraint, or more likely, the prompt instruction "Use the tag 'guess:'" meant I should output the word *as* the guess, but the previous turn's failure suggests the format might have been interpreted as needing the tag to be separate or the word to be the only content. Re{-}reading the specific error: "ABORT... violated the game rules". In the example provided in the prompt description, the format is "guess: apple". If I output "guess: tick", that matches the example. Perhaps the issue was that I didn't include an explanation? No, the prompt says "Use the tag 'guess:' to provide your guess" and "Do not generate any other text." It does not explicitly ask for an explanation tag like in Wordle, but the previous example (Example 2) had an explanation. Wait, the current prompt says "Use the tag 'guess:' to provide your guess" and "Do not generate any other text." This implies the output should *only* be the guess line. If I output "guess: tick", that is exactly what was done. Why was it aborted? Maybe the word "tick" is not 5 letters? "tick" is 4 letters. The rules state: "a valid English word of five lowercase letters". "tick" is 4 letters. That is the mistake. I need to ensure the word is exactly 5 letters long.} \\
\\ 
\texttt{New plan: I must ensure the guessed word is exactly five letters long. In the previous attempt, I guessed "tick" which is only four letters. I will select a five{-}letter word that fits the clue "tick off" (e.g., "check", "mark", "done", "clear", "finish") and output it strictly in the format "guess: {[}word{]}" with no extra text.} \\
            }
        }
    }
    & & \\ \\

\end{supertabular}
}

\end{document}
